%% toc: P15 165 165.10
\chapter{Lower Previsions Compared with Precise Probability, Accepted Belief and Decision Criteria}
\label{sec:p15-w02b-p15-036}

A coherent lower prevision fixes a closed convex set of precise models, but
it does not fix a choice. Several rules turn the set into a choice from a
menu of acts, and they need not agree. This chapter compares them in one
setting. Its inputs are a coherent lower prevision on a full space of
gambles and a finite menu of acts expressed as gambles on a linear utility
scale. Its conclusions are of three kinds. First, the rules are nested: the
optimistic and pessimistic scalar rules choose inside the set of maximal
acts, the acts that are best under some compatible model are maximal, and
every maximal act survives interval dominance
(\cref{thm:p15-w02b-p15-036-inclusions}). Second, all the rules coincide
with expected-utility maximization exactly when the lower prevision is
precise (\cref{thm:p15-w02b-p15-036-precise}), and a single menu can
separate all of them when it is not
(\cref{thm:p15-w02b-p15-036-menu}). Third, accepting an event licenses
deciding as if it were true exactly when its lower probability is one;
acceptance at a lower level, including the conjunction-closed acceptance of
a necessity measure, does not (\cref{thm:p15-w02b-p15-036-accepted,%
thm:p15-w02b-p15-036-threshold}). Under the vacuous model the rules reduce
to decision criteria that use only the set of possible states
(\cref{thm:p15-w02b-p15-036-vacuous}).

The gamble framework, coherence, the compatible set and the envelope
theorem are those of \cref{sec:p15-imprecise-probability}. The two-act
comparisons of maxmin preference with robust dominance are treated in
\cref{thm:p15-w02b-p15-032-ellsberg,thm:p17-robust-dominance}; this chapter
concerns menus and the relations among whole choice sets.

\section{Setting and five choice rules}
\label{sec:p15-w02b-p15-036-setting}

\begin{hbconvention}{def:p15-w02b-p15-036-standing}{Standing assumptions}
Let \(\Omega\), \(\Sigma\) and \(B=B(\Sigma)\) be as in
\cref{def:p15-ip-domain}, and let \(\underline P\) be a coherent lower
prevision on all of \(B\), with conjugate upper prevision
\(\overline P(f)=-\underline P(-f)\). Since the input is coherent on
\(B\), it equals its own natural extension
(\cref{thm:p15-ip-natural-extension}). Write
\(\mathcal M=\mathcal M(\underline P)\) for its compatible set of
\cref{def:p15-ip-compatible}. By
\cref{thm:p15-ip-representation,thm:p15-ip-dual}, \(\mathcal M\) is
nonempty, convex and weak-star compact, and
\[
 \underline P(f)=\min_{L\in\mathcal M}L(f),\qquad
 \overline P(f)=\max_{L\in\mathcal M}L(f)\qquad(f\in B),
\]
both extrema being attained. A \emph{menu} is a nonempty set
\(\mathcal A\subseteq B\) of acts; unless stated otherwise it is finite.
An act's values are utilities on one linear scale, so that sums and
differences of acts are meaningful and every compatible model values an
act by its expectation. The choice is made once, from the whole menu, with
no later information.
\end{hbconvention}

The linear-scale clause is a modelling assumption, not a consequence of
coherence. Miranda lists the linearity of the utility scale among the open
difficulties of the theory, and also notes that it is not settled which
optimality criterion suits which situation
\cite[§7, p.~38]{MirandaCoherentLowerPrevisionsCapture}.

\begin{hbdefinition}{def:p15-w02b-p15-036-rules}{Five choice rules for a lower prevision}
Under \cref{def:p15-w02b-p15-036-standing}, let \(\mathcal A\) be a menu
and \(f\in\mathcal A\).
\begin{enumerate}
\item For \(\alpha\in[0,1]\) put
 \(V_\alpha(f)=\alpha\,\underline P(f)+(1-\alpha)\,\overline P(f)\).
 The \emph{\(\alpha\)-maxmin choice set}
 \(\operatorname{ch}_\alpha(\mathcal A)\) is the set of maximizers of
 \(V_\alpha\) on \(\mathcal A\). The case \(\alpha=1\) is
 \emph{\(\Gamma\)-maximin}; the case \(\alpha=0\) maximizes the upper
 prevision.
\item \(f\) is \emph{interval undominated} if there is no
 \(g\in\mathcal A\) with \(\underline P(g)>\overline P(f)\). Write
 \(\operatorname{ch}_{\mathrm{id}}(\mathcal A)\) for these acts.
\item \(f\) is \emph{maximal} if there is no \(g\in\mathcal A\) with
 \(\underline P(g-f)>0\). Write
 \(\operatorname{ch}_{\mathrm{mx}}(\mathcal A)\) for these acts.
\item \(f\) is \emph{E-admissible} if some \(L\in\mathcal M\) satisfies
 \(L(f)\geq L(g)\) for every \(g\in\mathcal A\). Write
 \(\operatorname{ch}_{\mathrm E}(\mathcal A)\) for these acts.
\end{enumerate}
\end{hbdefinition}

The \(\alpha\)-weighted rule, its pessimism convention (\(\alpha=1\) is the
cautious extreme), the unanimity comparison, which Hill notes is called
maximality in statistics, and E-admissibility are presented by Hill for a
set of priors \cite[Sec.~2, p.~497]{Hill2021DecisionUncertainty}; the
finite-state versions are \cref{def:p17-rat84-prior-rules}. Miranda names
E-admissibility, interval dominance, \(\Gamma\)-maximinity and maximality
as generalizations of maximum expected utility to coherent lower
previsions, without stating their definitions
\cite[§6, p.~37]{MirandaCoherentLowerPrevisionsCapture}. The definitions
above are therefore fixed here, in the notation of
\cref{sec:p15-imprecise-probability}.

Two features of item~3 need to be kept in view. The test
\(\underline P(g-f)>0\) holds exactly when \(L(g)>L(f)\) for every
\(L\in\mathcal M\), because the minimum defining \(\underline P(g-f)\) is
attained; it is the strict desirability relation \(\succ_D^+\) of
\cref{def:p15-w02b-p15-032-rules}. It is not the strict part of Hill's
unanimity relation, which asks for a weak inequality under every model and
a strict one under some model; the undominated sets of the two relations
differ (\cref{thm:p15-w02b-p15-036-menu}, item~5). And items~2 and~3 are
pairwise tests, while item~4 asks for one model that ranks \(f\) first
against the entire menu.

\section{Inclusions and the precise case}
\label{sec:p15-w02b-p15-036-inclusions}

\begin{hbtheorem}{thm:p15-w02b-p15-036-inclusions}{proved=annex:p15.w02b-p15-036.inclusions}{The choice sets are nested}
Under \cref{def:p15-w02b-p15-036-standing}, let \(\mathcal A\) be a finite
menu.
\begin{enumerate}
\item Every set \(\operatorname{ch}_\alpha(\mathcal A)\),
 \(\operatorname{ch}_{\mathrm E}(\mathcal A)\),
 \(\operatorname{ch}_{\mathrm{mx}}(\mathcal A)\) and
 \(\operatorname{ch}_{\mathrm{id}}(\mathcal A)\) is nonempty.
\item \(\operatorname{ch}_\alpha(\mathcal A)\subseteq
 \operatorname{ch}_{\mathrm{mx}}(\mathcal A)\) for every
 \(\alpha\in[0,1]\); in particular every \(\Gamma\)-maximin act is maximal.
\item \(\operatorname{ch}_0(\mathcal A)\subseteq
 \operatorname{ch}_{\mathrm E}(\mathcal A)\).
\item \(\operatorname{ch}_{\mathrm E}(\mathcal A)\subseteq
 \operatorname{ch}_{\mathrm{mx}}(\mathcal A)\subseteq
 \operatorname{ch}_{\mathrm{id}}(\mathcal A)\).
\item If \(\mathcal A=\{f,g\}\) has two acts, then
 \(\operatorname{ch}_{\mathrm E}(\mathcal A)=
 \operatorname{ch}_{\mathrm{mx}}(\mathcal A)\), and both acts are maximal
 exactly when \(\underline P(f-g)\leq0\leq\overline P(f-g)\).
\item For an infinite menu every one of these sets can be empty.
\end{enumerate}
\end{hbtheorem}

Item~3 has no counterpart for \(\alpha>0\): the menu of
\cref{thm:p15-w02b-p15-036-menu} has a \(\Gamma\)-maximin act that is not
E-admissible. Item~5 says that the separation between maximality and
E-admissibility needs at least three acts. The attained maximum of
\cref{def:p15-w02b-p15-036-standing} is used in items~3 and~5; the remaining
items use only superadditivity and conjugacy.

\begin{hbtheorem}{thm:p15-w02b-p15-036-precise}{proved=annex:p15.w02b-p15-036.precise}{Agreement with expected utility exactly in the precise case}
Under \cref{def:p15-w02b-p15-036-standing}, the following are equivalent.
\begin{enumerate}
\item \(\underline P(f)=\overline P(f)\) for every \(f\in B\).
\item \(\mathcal M\) has exactly one element \(P\).
\item On every finite menu the sets \(\operatorname{ch}_\alpha\)
 (\(\alpha\in[0,1]\)), \(\operatorname{ch}_{\mathrm E}\),
 \(\operatorname{ch}_{\mathrm{mx}}\) and \(\operatorname{ch}_{\mathrm{id}}\)
 coincide.
\end{enumerate}
When they hold, every one of these sets equals the set of maximizers of
\(P(f)\) on the menu, and any two of its members have equal expectation.
When they fail, there are \(f\in B\) and a constant act \(c\) such that
\(\operatorname{ch}_1(\{f,c\})=\{c\}\) while
\(\operatorname{ch}_{\mathrm{mx}}(\{f,c\})=\{f,c\}\).
\end{hbtheorem}

In the precise case a tie means equal expectations, so the decision maker
is indifferent among the chosen acts. When the model is imprecise, the
maximal set can contain acts with different lower and upper previsions, as
\(f\) and \(c\) in the last clause; two maximal acts are then unranked, not
equally good. The same distinction appears for one-sided prices in
\cref{thm:p15-w02b-p15-032-unique}, where two-sided fair prices force one
distribution.

\section{One menu that separates the rules}
\label{sec:p15-w02b-p15-036-menu}

\begin{hbexample}{def:p15-w02b-p15-036-menu}{Four acts under an interval probability}
Let \(\Omega=\{\omega_1,\omega_2\}\) with \(\Sigma=\mathcal P(\Omega)\),
write an act as the pair \((f(\omega_1),f(\omega_2))\), and let \(P_t\)
be the probability with \(P_t(\{\omega_1\})=t\). Take
\[
 \mathcal M=\{P_t:\ 1/4\leq t\leq3/4\},
\]
with lower envelope \(\underline P\), and the menu
\(\mathcal A=\{a,b,c,d\}\) with
\[
 a=(1,0),\qquad b=(0,9/10),\qquad c=(2/5,2/5),\qquad d=(19/20,-1/20).
\]
Act \(a\) bets on \(\omega_1\), \(b\) bets on \(\omega_2\) at a smaller
stake, \(c\) is a sure payment, and \(d\) is \(a\) less a sure \(1/20\).
\end{hbexample}

\begin{hbtheorem}{thm:p15-w02b-p15-036-menu}{proved=annex:p15.w02b-p15-036.menu}{The rules on the four-act menu}
In \cref{def:p15-w02b-p15-036-menu}:
\begin{enumerate}
\item \(\underline P\) is coherent and \(\mathcal M(\underline P)=\mathcal M\).
 The intervals \([\underline P,\overline P]\) of \(a,b,c,d\) are
 \([1/4,3/4]\), \([9/40,27/40]\), \([2/5,2/5]\) and \([1/5,7/10]\).
\item \(\operatorname{ch}_\alpha(\mathcal A)=\{a\}\) for
 \(\alpha<7/10\), \(\{a,c\}\) for \(\alpha=7/10\), and \(\{c\}\) for
 \(\alpha>7/10\). In particular \(\operatorname{ch}_1(\mathcal A)=\{c\}\)
 and \(\operatorname{ch}_0(\mathcal A)=\{a\}\).
\item \(\operatorname{ch}_{\mathrm E}(\mathcal A)=\{a,b\}\):
 \(a\) is best under \(P_t\) for \(t\geq9/19\), \(b\) for \(t\leq9/19\),
 and \(c\), \(d\) under no \(P_t\).
\item \(\operatorname{ch}_{\mathrm{mx}}(\mathcal A)=\{a,b,c\}\) and
 \(\operatorname{ch}_{\mathrm{id}}(\mathcal A)=\{a,b,c,d\}\).
\item On the menu \(\{(0,0),(0,1)\}\) with the same \(\mathcal M\)
 enlarged to all \(P_t\), \(0\leq t\leq1\), both acts are E-admissible and
 maximal, but \((0,1)\) is weakly preferred under every model and strictly
 under \(P_0\), so \((0,0)\) is excluded by the strict part of unanimity.
\end{enumerate}
\end{hbtheorem}

So
\(\{a\}=\operatorname{ch}_0\subsetneq\operatorname{ch}_{\mathrm E}
\subsetneq\operatorname{ch}_{\mathrm{mx}}\subsetneq
\operatorname{ch}_{\mathrm{id}}\), and
\(\operatorname{ch}_1=\{c\}\) is maximal but disjoint from
\(\operatorname{ch}_{\mathrm E}\): every inclusion of
\cref{thm:p15-w02b-p15-036-inclusions} is strict on one menu, and
\(\Gamma\)-maximin selects a sure payment that no compatible model ranks
first. The precise comparison is item~3 read model by model: each \(P_t\)
chooses \(a\) or \(b\), and never \(c\). The act \(d\) shows what interval
dominance misses: \(a\) beats \(d\) by \(1/20\) in every state, yet their
intervals overlap, because the interval test compares the two acts under
two different models. The menu of \cref{thm:p17-rat84-criteria}, with
\(\mathcal M\) all of \(\Delta\), already separates E-admissibility from
\(\Gamma\)-maximin; the interval set used here keeps the separation when
the probability of \(\omega_1\) is bounded away from \(0\) and \(1\), and
adds the cases of \(\operatorname{ch}_0\) and interval dominance.

\section{Accepted events, thresholds and the vacuous model}
\label{sec:p15-w02b-p15-036-acceptance}

Qualitative belief accepts events rather than pricing gambles. The
question here is when acceptance of an event \(A\) licenses deciding as if
\(A\) were true, that is, comparing acts by their values on \(A\) alone.
Accepted events and the threshold postulate are those of
\cref{def:p15-r47-acceptance}; for a lower prevision the confidence in
\(A\) is its lower probability \(\underline P(A)\).

\begin{hbtheorem}{thm:p15-w02b-p15-036-accepted}{proved=annex:p15.w02b-p15-036.accepted}{Full acceptance and deciding as if}
Under \cref{def:p15-w02b-p15-036-standing}, let \(A\in\Sigma\). The
following are equivalent.
\begin{enumerate}
\item \(\underline P(A)=1\).
\item \(L(A)=1\) for every \(L\in\mathcal M\).
\item \(\underline P(f)=\underline P(g)\) whenever \(f,g\in B\) agree on
 \(A\).
\end{enumerate}
If they hold, then for every finite menu \(\mathcal A\) and every map
\(\varphi:\mathcal A\to B\) with \(\varphi(f)=f\) on \(A\) and
\(\varphi\) injective, \(\varphi\) carries each of the sets
\(\operatorname{ch}_\alpha(\mathcal A)\),
\(\operatorname{ch}_{\mathrm E}(\mathcal A)\),
\(\operatorname{ch}_{\mathrm{mx}}(\mathcal A)\),
\(\operatorname{ch}_{\mathrm{id}}(\mathcal A)\) onto the corresponding set
for \(\varphi(\mathcal A)\). If they fail, the acts \(0\) and
\(-I_{A^c}\) agree on \(A\) but
\(\operatorname{ch}_1(\{0,-I_{A^c}\})=\{0\}\).
\end{hbtheorem}

\begin{hbtheorem}{thm:p15-w02b-p15-036-threshold}{proved=annex:p15.w02b-p15-036.threshold}{A conjunction-closed acceptance that does not license deciding as if}
Let \(\Omega=\{\omega_1,\omega_2\}\), \(A=\{\omega_1\}\), and
\(\mathcal M=\{P_t:9/10\leq t\leq1\}\) in the notation of
\cref{def:p15-w02b-p15-036-menu}, with lower envelope \(\underline P\).
\begin{enumerate}
\item On events, \(\underline P\) is the necessity measure of the
 possibility distribution \(\pi(\omega_1)=1\), \(\pi(\omega_2)=1/10\):
 \(\underline P(A)=9/10\), \(\underline P(A^c)=0\). So \(A\) is accepted
 at every threshold \(\beta\in(0,9/10]\), and the events of positive
 lower probability are closed under intersection.
\item For \(f=(1,-100)\) and \(g=(0,0)\), comparing values on \(A\)
 strictly prefers \(f\). But \(\underline P(f-g)=-91/10\) and
 \(\overline P(f-g)=1\), so
 \(\operatorname{ch}_{\mathrm{mx}}(\{f,g\})=
 \operatorname{ch}_{\mathrm E}(\{f,g\})=\{f,g\}\), while
 \(\operatorname{ch}_1(\{f,g\})=\{g\}\).
\item Every precise \(P\) with \(P(A)=9/10\) gives \(P(f)=-91/10\) and
 chooses \(g\).
\end{enumerate}
\end{hbtheorem}

Accepting \(A\) in the conjunction-closed sense of
\cref{thm:p15-r47-threshold} is a statement about which events are
believed. It is not a licence to discard the states outside \(A\) when
acts differ there by large amounts: under every rule of
\cref{def:p15-w02b-p15-036-rules} the act that does best on \(A\) is
either not chosen or left unranked against the alternative. The
precise-probability comparison in item~3 makes the same point for a
Lockean threshold: \(P(A)=9/10\) supports belief in \(A\) at that
threshold, and expected utility still rejects acting on it. Only full
acceptance, \(\underline P(A)=1\), removes the states outside \(A\) from
every compatible model, and then by
\cref{thm:p15-w02b-p15-036-accepted} no rule can see them.

\begin{hbtheorem}{thm:p15-w02b-p15-036-vacuous}{proved=annex:p15.w02b-p15-036.vacuous}{The vacuous model and decision with only a set of possible states}
Let \(\Omega\) be finite, \(\Sigma=\mathcal P(\Omega)\), and
\(\underline P(f)=\min_{\omega\in\Omega}f(\omega)\). Then \(\underline P\)
is coherent, \(\mathcal M\) is the set of all probability mass functions on
\(\Omega\), and for a finite menu \(\mathcal A\):
\begin{enumerate}
\item \(V_\alpha(f)=\alpha\min f+(1-\alpha)\max f\); so
 \(\operatorname{ch}_1\) maximizes the worst outcome and
 \(\operatorname{ch}_0\) the best outcome;
\item \(f\) is maximal exactly when no \(g\in\mathcal A\) has
 \(g(\omega)>f(\omega)\) for every \(\omega\);
\item \(f\) is interval undominated exactly when no \(g\in\mathcal A\) has
 \(\min g>\max f\);
\item \(f\) is E-admissible exactly when some mass function \(p\) gives
 \[
  \textstyle\sum_\omega p(\omega)f(\omega)\geq\sum_\omega p(\omega)g(\omega)
  \qquad(g\in\mathcal A).
 \]
\end{enumerate}
\end{hbtheorem}

Item~1 identifies \(\Gamma\)-maximin and the \(\alpha=0\) rule under
complete ignorance with the pessimistic and optimistic criteria that use
only the worst and the best outcome. Dubois and Prade report that
possibility-theoretic counterparts of expected utility extend Wald's
pessimistic and optimistic criteria to possibility distributions
\cite[p.~308]{DuboisPrade2021RepresentingBelief}; that qualitative theory
and its axioms are not derived here. The two-state menu of
\cref{thm:p17-rat84-criteria} is the case \(\Omega=\{\omega_1,\omega_2\}\)
of \cref{thm:p15-w02b-p15-036-vacuous}: there maximin picks the sure
\((.4,.4)\), which item~4 excludes.

\section{Limits of the comparison}
\label{sec:p15-w02b-p15-036-limits}

\begin{hbwarning}{def:p15-w02b-p15-036-limits}{What the comparison does not settle}
\begin{enumerate}
\item None of the rules is derived from coherence. The model fixes
 \(\mathcal M\); the rule is a further input. The theorems above compare
 the rules once a model is given; they do not show that one rule is
 correct.
\item A maximal set with more than one act does not state indifference.
 Reporting it, deferring, or applying a further rule are separate choices;
 the \(\Gamma\)-maximin selection in \cref{thm:p15-w02b-p15-036-menu}
 shows that the further rule can select an act that no compatible model
 ranks first.
\item E-admissibility depends on the whole menu, so adding an act can
 remove an existing act from \(\operatorname{ch}_{\mathrm E}\) without
 changing any pairwise comparison between the existing acts; in
 \cref{def:p15-w02b-p15-036-menu}, \(c\) is E-admissible in \(\{a,c\}\)
 by \cref{thm:p15-w02b-p15-036-inclusions}, item~5, but not in
 \(\{a,b,c\}\).
\item The comparisons are for a single choice from a fixed menu. How a
 rule behaves when choices are made in stages after updating is not
 addressed; conditioning is treated in \cref{sec:p15-ip-conditioning} and
 \cref{sec:p15-w02b-p15-034-partition-updates}.
\item The utility scale is assumed linear and known. If utilities are
 themselves imprecise, the act values are not the gambles used here.
\item Finite menus are assumed; by
 \cref{thm:p15-w02b-p15-036-inclusions}, item~6, an infinite menu can
 leave every rule without a choice.
\end{enumerate}
\end{hbwarning}

\section{Relationships}
\label{sec:p15-w02b-p15-036-relations}

\paragraph{Precise probability.}
\Cref{thm:p15-w02b-p15-036-precise} is the decision form of the
uniqueness of fair prices in \cref{thm:p15-w02b-p15-032-unique}: two-sided
prices on all gambles force one model, and then every rule is
expected-utility maximization. The finite expected-utility representation
used for one precise model is \cref{sec:p17-finite-expected-utility}.

\paragraph{Two-act comparisons.}
Item~2 of \cref{thm:p15-w02b-p15-036-inclusions} is the menu form of the
implication from strict robust preference to a strict maxmin comparison in
\cref{thm:p15-w02b-p15-032-ellsberg,thm:p15-r47-price-boundaries,%
thm:p17-robust-dominance}. In the independence example of
\cref{thm:p15-w02b-p15-034-independence-decision}, the menu has two acts,
so by item~5 the unranked pair under epistemic independence is also a pair
of E-admissible acts.

\paragraph{Belief functions, possibility and ranks.}
Belief functions and necessity measures are lower envelopes of particular
compatible sets (\cref{thm:p15-r47-mobius,thm:p15-w02b-p15-034-strict-classes,%
thm:p15-ip-possibility-bridge}); every rule of
\cref{def:p15-w02b-p15-036-rules} applies to them through the gamble
extension of \cref{thm:p15-w02b-p15-034-belief-extension}. The example in
\cref{thm:p15-w02b-p15-036-threshold} is such a necessity measure. Ranking
functions correspond to consonant belief functions on a finite space
(\cref{thm:p15-w02b-p15-039-consonant}); the rules compared here use only a
lower envelope and say nothing about the ranking dynamics.

\paragraph{Decision under ambiguity.}
For finite state spaces the \(\alpha\)-maxmin and unanimity rules over a
prior set, with their representation theorems and the Ellsberg urns, are in
\cref{sec:p17-rat84-multiple-priors,sec:p17-ambiguity-and-robust-choice}.
The present chapter uses the compatible set of a coherent lower prevision
as the prior set; whether a decision maker's preferences are represented by
such a set is the separate question answered by those representation
theorems.

\paragraph{Sources.}
The definitions of \cref{def:p15-w02b-p15-036-rules} follow the names in
\cite[Sec.~2, p.~497]{Hill2021DecisionUncertainty} and
\cite[§6, p.~37]{MirandaCoherentLowerPrevisionsCapture}. The theorems and
examples of this chapter, with their proofs in the annex, are constructed
for it and are not taken from either source.

\index{lower prevision!decision rules}\index{E-admissibility}
\index{maximality}\index{interval dominance}\index{Gamma-maximin@\(\Gamma\)-maximin}
\index{accepted belief!and decision}
